\documentclass[11pt]{ctexart}
\usepackage{amsmath,amssymb,longtable}
\begin{document}
\section*{P35: 当前完整数学模型（2026-10-20）}
\section*{修订依据与边界}保留每例原4/64/192期数量及有限电池/功率域，补齐非零异质行驶耗能、4小时采样尺度、有限分时价格与终端70%电量目标；参数由本地情景给定。




\section*{Retained Core Model}
According to the \texttt{tex} formulas and the PDF, the directly solvable core model is the
linear programming reformulation on Page 20, denoted by $(R''_{\mathcal K})$. The original
continuous-time robust problem $(R)$ on Page 13 and the intermediate robust counterpart
$(R'_{\mathcal K})$ on Page 19 are omitted because they are intermediate non-convex or
semi-infinite formulations. The paper proves that the original robust V2G problem is exactly
equivalent to the linear program retained below.

\paragraph{Decision variables.}
For each period $k \in \mathcal{K}$, $x_k^b$ is the baseline charging power, $x_k^r$ is the
regulation bid, $m_k$ is an auxiliary epigraph variable, and $z$ is the epigraph variable for
the terminal piecewise-affine penalty. The additional variables
$\lambda_k^+,\lambda_k^-,\theta_k^+,\theta_k^-,\Lambda_{k,l}^+,\Lambda_{k,l}^-,
\Theta_{k,l}^+,\Theta_{k,l}^-$ are dual variables introduced by robust optimization.

\paragraph{Parameters and index sets.}
The feasible charging set is $\mathcal{X}_{\mathcal K}$. The bounds
$\bar y_k^+,\bar y_k^-,\bar y,\underline y,\bar y_0,\underline y_0,
\hat{\bar y}_0,\hat{\underline y}_0$ are battery and initial-state limits.
The parameters $\eta^+,\eta^-,\Delta t,\gamma,\hat{\gamma},d_k,r_n,q_n$ are as in the paper,
with $\Delta\eta = \frac{1}{\eta^-} - \eta^+ > 0$.
The terminal penalty is represented by breakpoint sets
$\mathcal{N}_+,\mathcal{N}_-,\mathcal{N}_0$.
The horizon indices are
$I(k,l)=\min\{k,\ l+\Gamma/\Delta t-1\}$ and
$\hat I(K,k)=\min\{K,\ k+\hat{\Gamma}/\Delta t-1\}$.

\section*{Linear Programming Reformulation}
\begin{equation}
\min \quad c_{\mathcal K}(\boldsymbol{x}^b,\boldsymbol{x}^r) + z
\tag{R''_{\mathcal K}-Obj}
\end{equation}

\begin{equation}
\boldsymbol{x}^b,\boldsymbol{x}^r \in \mathcal{X}_{\mathcal K},\quad
z \in \mathbb{R},\quad
\boldsymbol{m},\boldsymbol{\lambda}^+,\boldsymbol{\lambda}^-,
\boldsymbol{\theta}^+,\boldsymbol{\theta}^- \in \mathbb{R}_+^K,
\tag{LP-a}
\end{equation}

\begin{equation}
\boldsymbol{\Lambda}^+,\boldsymbol{\Lambda}^-,
\boldsymbol{\Theta}^+,\boldsymbol{\Theta}^- \in \mathbb{R}_+^{K \times K}
\tag{LP-b}
\end{equation}

\begin{equation}
x_k^r + x_k^b \le \bar y_k^+,\qquad
x_k^r - x_k^b \le \bar y_k^-
\quad \forall k \in \mathcal K
\tag{27}
\end{equation}

\begin{equation}
m_k \ge \eta^+ x_k^r,\qquad
m_k \ge \frac{1}{\eta^-}x_k^r - \Delta\eta\,x_k^b
\quad \forall k \in \mathcal K
\tag{28}
\end{equation}

\begin{equation}
\sum_{l=1}^{k} \Delta t \left( \eta^+x_l^b + \Lambda_{k,l}^+ - d_l \right)
\;+\; \gamma \sum_{l=1}^{k} \Theta_{k,l}^+
\le \bar y - \bar y_0
\quad \forall k \in \mathcal K \cup \{0\}
\tag{29}
\end{equation}

\begin{equation}
\sum_{l=1}^{k} \Delta t \left( \eta^+x_l^b - \Lambda_{k,l}^- - d_l \right)
\;-\; \gamma \sum_{l=1}^{k} \Theta_{k,l}^-
\ge \underline y - \underline y_0
\quad \forall k \in \mathcal K \cup \{0\}
\tag{30}
\end{equation}

\begin{equation}
\sum_{k \in \mathcal K} \Delta t \left( \eta^+x_k^b + \lambda_k^+ - d_k \right)
\;+\; \hat{\gamma}\sum_{k \in \mathcal K}\theta_k^+
\le \frac{z - r_n}{q_n} - \hat{\bar y}_0
\quad \forall n \in \mathcal N_+
\tag{31}
\end{equation}

\begin{equation}
\sum_{k \in \mathcal K} \Delta t \left( \eta^+x_k^b - \lambda_k^- - d_k \right)
\;-\; \hat{\gamma}\sum_{k \in \mathcal K}\theta_k^-
\ge \frac{z - r_n}{q_n} - \hat{\underline y}_0
\quad \forall n \in \mathcal N_-
\tag{32}
\end{equation}

\begin{equation}
r_n \le z
\quad \forall n \in \mathcal N_0
\tag{33}
\end{equation}

\begin{equation}
\Lambda_{k,l}^+ + \sum_{i=l}^{I(k,l)} \Theta_{k,i}^+ \ge \eta^+ x_l^r,\qquad
\Lambda_{k,l}^- + \sum_{i=l}^{I(k,l)} \Theta_{k,i}^- \ge m_l
\quad \forall k,l \in \mathcal K:\ l \le k
\tag{34}
\end{equation}

\begin{equation}
\lambda_k^+ + \sum_{i=k}^{\hat I(K,k)} \theta_i^+ \ge \eta^+ x_k^r,\qquad
\lambda_k^- + \sum_{i=k}^{\hat I(K,k)} \theta_i^- \ge m_k
\quad \forall k \in \mathcal K
\tag{35}
\end{equation}


\section*{当前非退化运营数据}
每个实例保留原时段数、有限电池及功率域。$d_k>0$为异质行驶耗能率，终端分段惩罚
为$\max\{0,.6(y-0.7\bar y),.6(0.7\bar y-y)\}$；效率$\eta^+=\eta^-=.85$。
时段为4小时运营窗的等长采样，具体分时价格、初始电量及合法不确定集窗口由实例明确给定。
原零需求案例作为历史参考数据保存。


\end{document}
